% TLDR: Prospective counts separate artifact instances, identities, and transfer claims.
\section{Prospective Study Design}
\label{sec:design}
The starting population is the fixed public VeriBench release \texttt{veribench@1.3}, containing 896 Python tasks.
The data source is pinned to commit \texttt{c9dfd8e} of the public Harbor dataset collection \citep{veribenchrelease}.
Task contents, reference versions, instruction version, and toolchain must be hashed before collection.
The four required solver configurations are Claude Code, Codex, a Gemini agent, and a Grok agent.
These names define the intended panel; exact model versions, harnesses, tools, effort, and budgets remain to be frozen.
They are deliberately selected configurations, not random draws from every possible future solver.

\subsection{Task and human allocation}
The initial allocation is 64 training tasks, 64 validation tasks, and 768 held-out tasks.
Related task variants and near-duplicates are grouped before splitting.
If those groups prevent the exact counts, the allocation and budget must be amended before test inspection.
We recruit 24 qualified experts: eight for training and validation, and 16 different raters for the final audit.
Each artifact receives three independently submitted ratings from distinct randomly assigned raters in the relevant pool.
Assignment balances workload and maintains overlap between raters; it is recorded rather than inferred from a dense tensor.

The final audit covers all 896 task identities with fresh solver generations.
Consequently, its 128 previously seen tasks answer a fresh-output question, while its 768 held-out tasks answer a task-transfer question.
Those strata are reported separately even when a full-benchmark mean is also useful.
One solver output per task--configuration pair is collected initially in each stage, without selecting the best of retries.
This design does not estimate generation variance at a fixed task--configuration pair.
Repeated identical bytes remain distinct generation records but do not create evidence of independent content.

We randomly select 128 of the held-out tasks before observing their ratings and collect eight total ratings for each of their four artifacts.
These additional five ratings support distribution diagnostics; they do not make the empirical category histogram exact.
Table~\ref{tab:budget} distinguishes new artifacts from additional ratings of an existing artifact.
The resulting workload is 192 ratings per training/validation rater and 832 per audit rater under exact balancing.
These are planned counts, not measured or assumed feasible time commitments.

\begin{table}[t]
\centering
\begin{tabular}{lrr}
\toprule
Collection stage & Artifacts & Ratings \\
\midrule
128 tasks $\times$ 4, train/validation & 512 & 1,536 \\
896 tasks $\times$ 4, final audit & 3,584 & 10,752 \\
Five extra on 128 $\times$ 4 artifacts & 0 & 2,560 \\
\midrule
Total & 4,096 & 14,848 \\
\bottomrule
\end{tabular}
\caption{Prospective collection budget, not collected data.
The two stages have separate generations and disjoint human pools.
Three ratings are assigned initially to every artifact; the dense subset has eight in total.}
\label{tab:budget}
\end{table}

\subsection{What is held out}
The main judge is trained and selected with all four configurations' training and validation outputs.
Its audit establishes performance on the evaluated configurations, with separate new-task and new-human evidence.
It cannot also establish that these same configurations were unseen by that judge.

To test transfer to each excluded configuration, we fit four additional judges.
For each fit, one configuration's outputs, ratings, and prompt examples are excluded from both training and validation.
The fit is evaluated on that configuration's 768 held-out tasks using audit raters.
All four fits reuse the same collected audit ratings.
These experiments test the transfer of the training/selection procedure, not a single judge that somehow treats all four training configurations as unseen.
They also do not establish transfer to arbitrary future model versions.
Historical task and solver exposure is recorded; a new split cannot erase earlier prompt design or unknown pretraining exposure.

The final human pool supplies no ratings for judge selection.
An analysis conditional on the 16 recruited audit experts can establish agreement with that panel.
An inference about a wider expert population additionally requires a defensible recruitment and sampling model.
Holding out identities reduces direct leakage but does not by itself make either experts or solver configurations representative.

\begin{table}[t]
\centering
\begin{tabular}{p{0.35\columnwidth}p{0.55\columnwidth}}
\toprule
Audit stratum & Supported question \\
\midrule
Full audit (896 tasks) & Agreement on the complete fixed benchmark and four evaluated solver configurations. \\
128 seen tasks, fresh outputs & Agreement on new generations from tasks used for training or selection. \\
768 held-out tasks & Transfer across task groups untouched in the new selection procedure. \\
Four excluded-solver fits & Transfer of that procedure to each configuration excluded from its fit. \\
16 audit raters & Agreement with new raters, with wider population inference conditional on recruitment assumptions. \\
\bottomrule
\end{tabular}
\caption{The holdouts answer different questions.
No row alone establishes accuracy on every future task, solver configuration, or rater.}
\label{tab:holdouts}
\end{table}

\subsection{Pilot and fixed execution}
A training/validation-only pilot measures rating time, disagreement, fatigue, failure rates, and variance of paired losses.
It also traces one artifact through generation, extraction, human display, judge input, probability parsing, and scoring using content hashes.
As an artificial pipeline control, the pilot supplies a known rating in the prompt and checks that the returned category or probability vector reproduces it through parsing and normalization.
This rating-reproduction check diagnoses instruction and transport failures; it is never scored as evidence of semantic ability or included among test artifacts.
Before final ratings are opened, we freeze the rubric, assignments, judge versions, analysis, repeat count, failure policy, endpoints, and practical tolerances.
Task and rater exclusions, parser errors, missing outputs, and timeouts remain in the accounting.
The plan is not a completed preregistration: practical thresholds and the feasibility of the proposed workload remain unresolved.
